\section{Conclusion}

This seminar paper presented the path algebra for graph query languages proposed by Angles et al. \cite{angles2024path}. The algebra treats paths as first-class query objects and provides operators for selection, union, concatenation, and recursion. These operators make it possible to describe path queries in a compositional way.

The paper also examined how the algebra relates to path querying features in GQL and SQL/PGQ, including different path selection and restriction semantics. Furthermore, it was compared with relational algebra, automata-based approaches to regular path queries, and semiring-based models. These comparisons show that the path algebra combines ideas from existing formalisms while focusing specifically on the construction and manipulation of paths.

The main strength of the approach is that it provides a common logical framework for describing different types of path queries. At the same time, the algebra does not specify physical execution strategies or cost models. Efficient evaluation of recursive path queries, especially on large and cyclic graphs, therefore remains an important practical challenge.

Overall, the path algebra provides a useful foundation for reasoning about graph path queries and offers a clear connection between the semantics of modern graph query languages and their underlying algebraic operations.